% Part: lambda-calculus
% Chapter: introduction
% Section: church-rosser

\documentclass[../../../include/open-logic-section]{subfiles}

\begin{document}

\olfileid{lam}{int}{cr}
\olsection{ચર્ચ--રોસર ગુણધર્મ}

\begin{thm}
\ollabel{thm:church-rosser}
$M$, $N_1$ અને~$N_2$ એવાં પદો હોય કે $M \red N_1$ અને $M \red
N_2$. તો એવું પદ~$P$ છે કે $N_1 \red P$ અને $N_2 \red P$.
\end{thm}

\begin{cor}
માનો કે~$M$નું સામાન્ય સ્વરૂપમાં લઘુકરણ થઈ શકે છે. તો આ
સામાન્ય સ્વરૂપ અનન્ય છે.
\end{cor}

\begin{proof}
જો $M \red N_1$ અને $M \red N_2$, તો અગાઉના પ્રમેય મુજબ એવું
પદ~$P$ છે કે~$N_1$ અને~$N_2$ બંનેનું~$P$માં લઘુકરણ થાય છે.
$N_1$ અને~$N_2$ બંને સામાન્ય સ્વરૂપમાં હોય, તો આ ત્યારે જ બની
શકે જ્યારે $N_1 \ident P \ident N_2$.
\end{proof}

છેલ્લે, બે પદો~$M$ અને~$N$ કોઈ એક જ પદમાં લઘુકૃત થાય તો તેમને
\emph{$\beta$-સામ્યવાળાં}, અથવા ટૂંકમાં \emph{સામ્યવાળાં},
કહીશું. એટલે કે એવું~$P$ હોય કે $M \red P$ અને $N \red P$.
તેને $M \equal[\beta] N$ લખીએ છીએ. 
\olref{thm:church-rosser}નો ઉપયોગ કરીને ચકાસી શકાય કે
$\equal[\beta]$ સામ્ય-સંબંધ છે અને તેમાં વધારાનો ગુણધર્મ છે:
દરેક $M$ અને~$N$ માટે $M \red N$ અથવા $N \red M$ હોય તો
$M \equal[\beta] N$. (ખરેખર, આ ગુણધર્મ ધરાવતા સામ્ય-સંબંધોમાં
$\equal[\beta]$ \emph{સૌથી નાનો} છે એમ બતાવી શકાય.)

\end{document}
